% Source: https://artofproblemsolving.com/community/c3h2718639_circle_and_geometry_100th_post
% Stable identity: aops-topic-2718639-q01
Triangle $APQ$ is isosceles with $AP=AQ=1$ and apex angle $30^\circ$, so $PQ=2\sin15^\circ$. The circle is tangent to $PQ$ and to the extensions of $AP$ and $AQ$, so it is the excircle of triangle $APQ$ opposite $A$, and the tangent length from $A$ equals the semiperimeter:
\[AB=\dfrac{AP+AQ+PQ}{2}=1+\sin15^\circ.\]

\begin{asy}
usepackage("amsmath");
size(5cm);
defaultpen(fontsize(10pt));
real s=Sin(15),c=Cos(15);real AB=1+s;real r=AB*Tan(15);
pair A=(0,0),O=(AB/c,0),B=AB*dir(-15),C=AB*dir(15),P=dir(-15),Q=dir(15);
draw(A--1.9*dir(-15),linewidth(0.8),Arrow(4));draw(A--1.9*dir(15),linewidth(0.8),Arrow(4));
draw(circle(O,r));draw(B--C);draw(P--Q);dot(O);
label("$A$",A,W);label("$B$",B,S);label("$C$",C,N);label("$P$",P,S);label("$Q$",Q,N);label("$O$",O,E);
label("$30^\circ$",A+(0.42,0),fontsize(7pt));label("$1$",P/2,S,fontsize(7pt));
\end{asy}

\textbf{(a)} $OB\perp AB$ and $\angle OAB=15^\circ$, so
\[r=AB\tan15^\circ=(1+\sin15^\circ)\tan15^\circ=\boxed{\dfrac{8-4\sqrt3+3\sqrt6-5\sqrt2}{4}}\approx0.337,\]
using $\sin15^\circ=\dfrac{\sqrt6-\sqrt2}{4}$ and $\tan15^\circ=2-\sqrt3$.

\textbf{(b)} $BC=2AB\sin15^\circ=2\sin15^\circ+2\sin^215^\circ=\dfrac{\sqrt6-\sqrt2}{2}+1-\dfrac{\sqrt3}{2}=\boxed{\dfrac{2+\sqrt6-\sqrt2-\sqrt3}{2}}\approx0.652$.
